\section{Related Work}
\label{sec:related}

\subsection{Gaze-Contingent and Adaptive Reading}

Using the reader's gaze to change the reading experience in real time has a long line of demonstrations.
The eyeBook coupled an eye tracker to a digital book to trigger in-situ effects as the reader progressed \cite{biedert2010eyebook}.
More recently, \citet{medan2024newline} presented text as a single line that scrolls with the reader's webcam-estimated gaze and emphasises the fixated word; a pilot with 13 participants found no significant change in comprehension or speed but a clear preference for the interface.
\citet{rummens2026highlighting} coloured the word under fixation in real time for seventy second graders reading aloud and reported faster reading with shorter fixations, fewer regressions, and less re-reading, without loss of pronunciation accuracy or comprehension, although the children preferred the static text.
That split between performance and preference is a useful reporting model for adaptive reading.
Gaze during reading is also used as a measurement instrument in its own right, for example to compare how readers process AI-generated and human-authored text \cite{ilyas2025mind}.

Within the Reading the Reader programme, three master's projects built successive adaptive readers: a flexible reading and gaze-collection application \cite{refsgaard2024rtr}, an adaptive e-reader used to compare four intervention presentation forms, which were found to differ substantially in disruption and acceptance \cite{pereira2024typography}, and a reactive frontend designed to consume machine-learning output \cite{palinec2024reactive}.
\citet{jensen2025context} then showed that context preservation is where intervention design bites: in a 22-participant within-subjects study of four intervention designs they introduced reading-resume time as a recovery measure and reported that a gradual transition combined with context highlighting resumed fastest, reducing the time lost to an intervention from 14.9\,s to 4.0\,s, and was the most preferred design.
Each of these systems, however, was built for its own study.
Sensing, decision logic, and presentation were coupled in one application, so the next study copied and modified the previous system rather than reconfiguring a shared one.
Our platform is the infrastructural complement to this line of work: it is built so that exactly such comparisons of intervention designs can be run, varied, and reconstructed without touching core code.

\subsection{Real-Time Gaze Processing for Reading}

Reducing a raw gaze stream to oculomotor events is the classical first step, with identification algorithms catalogued by velocity, dispersion, and area-of-interest families \cite{salvucci2000identifying} and reading behaviour extensively characterised \cite{rayner1998eyemovements}.
Real-time use sharpens the problem: noisy vertical gaze must be assigned to the correct line while the reader is still reading.
\citet{kaltenberger2026line} address this with CONF-LA, a confidence-scored online fixation-to-line assignment that defers uncertain assignments.
Our contribution on this front is architectural rather than algorithmic.
The platform's built-in analysis is a dwell-based, token-level area-of-interest method whose spatial half hit-tests gaze against the words' live on-screen boxes (Section~\ref{sec:platform-mapping}), so a typographic reflow cannot invalidate the mapping; and the whole analysis stage is a replaceable, logged strategy, so a stronger detector or a confidence-aware line-assignment method can be attached as a provider without modifying the platform.
We do not claim token-level assignment accuracy for the built-in heuristic, which has not been validated against ground truth.

\subsection{Experiment Tooling}

A researcher building a reading study today would more likely reach for general-purpose experiment software than for a research prototype.
Tobii Pro Lab \cite{tobii2025prolab}, PsychoPy \cite{peirce2019psychopy}, OpenSesame \cite{mathot2012opensesame}, and Psychtoolbox \cite{brainard1997psychtoolbox} are mature platforms for experiment design, stimulus delivery, and gaze capture, and PsychoPy's ioHub already offers a common interface across tracker vendors together with a simulated mouse-driven gaze source \cite{psychopy2026iohub}, so hardware abstraction and simulated gaze are not in themselves new.
What these tools leave to study-specific programming is the reading-specific loop: mapping gaze onto reflowable text, proposing and gating typographic changes, committing them without losing the reader, and measuring what the change cost.
Table~\ref{tab:capability} summarises the authors' assessment of the three most established tools against the capabilities a gaze-driven reading experiment requires; none closes the loop from live gaze to a typographic change in the text being read, and none exposes a pluggable decision boundary.
Our platform occupies that region as a focused instrument for one family of experiments; it complements rather than replaces the general-purpose tools.

% TODO(authors): re-verify the Partial/No ratings against current primary documentation before submission.
\begin{widetable}
\centering
\small
\caption{Capability comparison of established experiment software with our platform, across the capabilities a gaze-driven reading experiment requires.
Entries reflect the authors' assessment of documented capabilities \cite{tobii2025prolab, peirce2019psychopy, brainard1997psychtoolbox}.}
\label{tab:capability}
\begin{tabular}{@{}lcccc@{}}
\toprule
Capability & Tobii Pro Lab & PsychoPy & Psychtoolbox & This platform \\
\midrule
Design experiments & Yes & Yes & Yes & Yes \\
Record gaze data & Yes & Partial & Partial & Yes \\
Analyse recorded data & Yes & Partial & No & Yes \\
Real-time gaze-driven text intervention & No & Partial & Partial & Yes \\
Pluggable decision modules & No & No & No & Yes \\
\bottomrule
\end{tabular}
\end{widetable}
